\section{Lista de verificación del proceso y validación}
\subsection{Preparación previa al entrenamiento}
Antes de entrar al entrenamiento self-rewarding, las recomendaciones de Weston incluyen: elegir las seeds de Self-Instruct, preparar los evaluators de SelfT, configurar el inference compute budget y establecer la frecuencia de cross-check. Estos pasos permiten que cada iteración genere de manera controlada una verified response, en la que «las tareas producidas por Self-Instruct deben cubrir matemáticas, programación y juicio» (SRT 674-676).

\begin{itemize}
    \item Establecer la seed de Self-Instruct (Math / Code / Reasoning / Dialogue) y ampliarla hasta obtener diverse prompts.
    \item Preparar el pipeline de verified response y definir los criterios de etiquetado de cross-check.
    \item Preentrenar o hacer fine-tuning del SelfT evaluator, para garantizar que el judge pueda distinguir entre draft y verified.
    \item Planificar el inference-time compute, para evitar que la salida de cross-check se convierta en un compute bottleneck (SRT 3236-3252).
    \item Respaldar los datos de conversación de Self-Feeding, para usarlos en la supervisión posterior.
\end{itemize}

\subsection{Procesamiento posterior a la evaluación}
Al terminar cada ronda de entrenamiento, es necesario usar los results para calibrar el meta judge, calcular el win rate y registrar si la verified response es consistente entre iteration. Weston menciona en particular el proceso de ``kept verifying itself and iterating'' (SRT 1908-2027), y recomienda extraer la verified response para compararla con Alpaca Eval 2, y volver a puntuarla con el judge, para asegurar que se rompa el plateau.

\begin{itemize}
    \item Hacer que el meta judge vuelva a puntuar draft y verified response, para capturar el drift en tiempo real.
    \item Comparar la verified response más reciente con Alpaca Eval 2 / Alpaca Eval 2 baseline.
    \item Reunir los error logs, prestando especial atención a las muestras de hallucination y overthinking.
    \item Ajustar el peso del loss del reward model y la frecuencia de cross-check según la retroalimentación del judge.
    \item Marcar las verified response de alta calidad como future seed, para seguir ampliando el corpus de Self-Instruct.
\end{itemize}

\subsection{Resumen de la sección}
La preparación previa al entrenamiento se encarga de garantizar que las distintas tareas y el judge estén listos, mientras que el procesamiento posterior a la evaluación se encarga de la corrección y la síntesis: solo si la verified response se usa tanto para el entrenamiento como para la evaluación, el self-rewarding puede seguir impulsando la calidad del reasoning.

